\pdfvariable suppressoptionalinfo 512\relax
\pdfvariable trailerid{}
\documentclass[10pt]{article}
\usepackage[a4paper,margin=0.68in]{geometry}
\usepackage{fontspec}
\usepackage{amsmath,amssymb,booktabs}
\newfontfamily\cjkfont{Droid Sans Fallback}
\linespread{0.96}
\title{\vspace{-1.5em}Every-Iterate Quadratic Amplitudes for Finite Cyclic Rotations}
\author{Route-A structural certificate C161}
\date{}
\begin{document}
\maketitle
\begin{abstract}
For the odd cyclic rotation $R_q(x)=x+1$ with quadratic observable
$ax^2+bx$, we evaluate the complete Birkhoff amplitude at every iterate and
for every source parameter.  If $d=(an,q)$, the sum vanishes exactly when
$d\nmid an(n-1)+bn$; otherwise its nonconstant branch has magnitude
$d\sqrt{q/d}$ and an explicit Jacobi sign and completed-square phase.  Over a
prime source ring the zero level is $1+(\Delta/p)$ in the quadratic branch;
for $p\ge5$ and $a=1,b=0$, $\Delta=n^2(1-n^2)/3$.  This all-parameter
theorem replaces a
Heisenberg candidate that failed its quotient-coordinate and local-equivalence
proof gate.  Exhaustive checks through odd $q\le31$ are regression sentinels,
not the proof.  A same-clock finite unitary realizes the amplitude as
$\operatorname{Tr}(U^nK^{-n})$.  No target trace law, arithmetic Euler data,
or Hilbert--P\'olya operator is claimed.
\end{abstract}
\noindent\textbf{Keywords:} cyclic dynamics; Birkhoff sum; quadratic Gauss
sum; exact vanishing; discriminant; finite Koopman unitary.

\medskip
\noindent{\cjkfont\parbox{\linewidth}{\raggedright\emergencystretch=2em
中\hspace{0pt}文\hspace{0pt}摘\hspace{0pt}要\hspace{0pt}：对\hspace{0pt}奇\hspace{0pt}数\hspace{0pt}阶\hspace{0pt}循\hspace{0pt}环\hspace{0pt}旋\hspace{0pt}转\hspace{0pt} $R_q(x)=x+1$ 及\hspace{0pt}二\hspace{0pt}次\hspace{0pt}观\hspace{0pt}测\hspace{0pt}量\hspace{0pt} $ax^2+bx$，\allowbreak
本\hspace{0pt}文\hspace{0pt}在\hspace{0pt}任\hspace{0pt}意\hspace{0pt}参\hspace{0pt}数\hspace{0pt}与\hspace{0pt}任\hspace{0pt}意\hspace{0pt}迭\hspace{0pt}代\hspace{0pt}下\hspace{0pt}精\hspace{0pt}确\hspace{0pt}求\hspace{0pt}出\hspace{0pt}完\hspace{0pt}整\hspace{0pt}比\hspace{0pt}尔\hspace{0pt}科\hspace{0pt}夫\hspace{0pt}振\hspace{0pt}幅\hspace{0pt}。\hspace{0pt}\allowbreak
令\hspace{0pt} $d=(an,q)$；\allowbreak 当\hspace{0pt}且\hspace{0pt}仅\hspace{0pt}当\hspace{0pt} $d$ 不\hspace{0pt}整\hspace{0pt}除\hspace{0pt} $an(n-1)+bn$ 时\hspace{0pt}振\hspace{0pt}幅\hspace{0pt}为\hspace{0pt}零\hspace{0pt}；\allowbreak
否\hspace{0pt}则\hspace{0pt}非\hspace{0pt}平\hspace{0pt}凡\hspace{0pt}分\hspace{0pt}支\hspace{0pt}的\hspace{0pt}模\hspace{0pt}长\hspace{0pt}为\hspace{0pt} $d\sqrt{q/d}$，\allowbreak 符\hspace{0pt}号\hspace{0pt}与\hspace{0pt}相\hspace{0pt}位\hspace{0pt}分\hspace{0pt}别\hspace{0pt}由\hspace{0pt}雅\hspace{0pt}可\hspace{0pt}比\hspace{0pt}符\hspace{0pt}号\hspace{0pt}和\hspace{0pt}
配\hspace{0pt}方\hspace{0pt}余\hspace{0pt}数\hspace{0pt}确\hspace{0pt}定\hspace{0pt}。\hspace{0pt}\allowbreak 素\hspace{0pt}数\hspace{0pt}阶\hspace{0pt}情\hspace{0pt}形\hspace{0pt}的\hspace{0pt}零\hspace{0pt}水\hspace{0pt}平\hspace{0pt}数\hspace{0pt}由\hspace{0pt}判\hspace{0pt}别\hspace{0pt}式\hspace{0pt} $\Delta$ 控\hspace{0pt}制\hspace{0pt}；\allowbreak
当\hspace{0pt} $p\ge5$ 时\hspace{0pt}，纯\hspace{0pt}二\hspace{0pt}次\hspace{0pt}子\hspace{0pt}族\hspace{0pt}满\hspace{0pt}足\hspace{0pt} $\Delta=n^2(1-n^2)/3$。\hspace{0pt}\allowbreak
此\hspace{0pt}前\hspace{0pt}海\hspace{0pt}森\hspace{0pt}堡\hspace{0pt}模\hspace{0pt}型\hspace{0pt}的\hspace{0pt}商\hspace{0pt}坐\hspace{0pt}标\hspace{0pt}及\hspace{0pt} $2$、\hspace{0pt}$5$ 进\hspace{0pt}局\hspace{0pt}部\hspace{0pt}等\hspace{0pt}价\hspace{0pt}未\hspace{0pt}完\hspace{0pt}成\hspace{0pt}严\hspace{0pt}格\hspace{0pt}证\hspace{0pt}明\hspace{0pt}，\allowbreak
故\hspace{0pt}按\hspace{0pt}硬\hspace{0pt}门\hspace{0pt}槛\hspace{0pt}更\hspace{0pt}换\hspace{0pt}模\hspace{0pt}型\hspace{0pt}。\hspace{0pt}\allowbreak 有\hspace{0pt}限\hspace{0pt}范\hspace{0pt}围\hspace{0pt}穷\hspace{0pt}举\hspace{0pt}仅\hspace{0pt}用\hspace{0pt}于\hspace{0pt}回\hspace{0pt}归\hspace{0pt}验\hspace{0pt}证\hspace{0pt}；\allowbreak
理\hspace{0pt}论\hspace{0pt}证\hspace{0pt}明\hspace{0pt}来\hspace{0pt}自\hspace{0pt}公\hspace{0pt}因\hspace{0pt}数\hspace{0pt}约\hspace{0pt}化\hspace{0pt}和\hspace{0pt}高\hspace{0pt}斯\hspace{0pt}和\hspace{0pt}求\hspace{0pt}值\hspace{0pt}。\hspace{0pt}\allowbreak
本\hspace{0pt}文\hspace{0pt}不\hspace{0pt}主\hspace{0pt}张\hspace{0pt}目\hspace{0pt}标\hspace{0pt}迹\hspace{0pt}公\hspace{0pt}式\hspace{0pt}、\hspace{0pt}欧\hspace{0pt}拉\hspace{0pt}数\hspace{0pt}据\hspace{0pt}或\hspace{0pt}希\hspace{0pt}尔\hspace{0pt}伯\hspace{0pt}特\hspace{0pt}波\hspace{0pt}利\hspace{0pt}亚\hspace{0pt}算\hspace{0pt}子\hspace{0pt}。\hspace{0pt}\par\smallskip
关\hspace{0pt}键\hspace{0pt}词\hspace{0pt}：循\hspace{0pt}环\hspace{0pt}动\hspace{0pt}力\hspace{0pt}学\hspace{0pt}；\allowbreak 比\hspace{0pt}尔\hspace{0pt}科\hspace{0pt}夫\hspace{0pt}和\hspace{0pt}；\allowbreak 二\hspace{0pt}次\hspace{0pt}高\hspace{0pt}斯\hspace{0pt}和\hspace{0pt}；\allowbreak
精\hspace{0pt}确\hspace{0pt}消\hspace{0pt}失\hspace{0pt}；\allowbreak 判\hspace{0pt}别\hspace{0pt}式\hspace{0pt}；\allowbreak 有\hspace{0pt}限\hspace{0pt}库\hspace{0pt}普\hspace{0pt}曼\hspace{0pt}酉\hspace{0pt}算\hspace{0pt}子\hspace{0pt}。\hspace{0pt}
}}

\section{All-iterate evaluation}
Let $q$ be odd, let $n\ge1$, let $R_q(x)=x+1$ on $\mathbb Z/q\mathbb Z$, and
$\phi_{a,b}(x)=ax^2+bx$.  Summing $1,j,j^2$ gives
\begin{align}
S_n\phi(x)&=A_nx^2+B_nx+C_n\pmod q,\tag{1}\\
A_n&=an,\quad B_n=an(n-1)+bn,\nonumber\\
C_n&=\frac{an(n-1)(2n-1)}6+\frac{bn(n-1)}2.\nonumber
\end{align}
The quotients in $C_n$ are integers before reduction.  Define
\[
G_{q,n}(a,b)=\sum_{x\bmod q}\exp(2\pi iS_n\phi(x)/q).
\]
Put $d=(A_n,q)$ and $Q=q/d$.  Translation $x\mapsto x+Q$ multiplies the
sum by $\exp(2\pi iB_n/d)$.  Hence $G_{q,n}=0$ exactly when $d\nmid B_n$.
If $d\mid B_n$ and $Q>1$, write $A'=A_n/d$, $B'=B_n/d$, and
\[
h\equiv-(4A')^{-1}(B')^2\pmod Q,\qquad r\equiv C_n+dh\pmod q.
\]
Reduction gives $d$ copies of a primitive odd quadratic sum, and completing
the square gives
\[
G_{q,n}(a,b)=d\left(\frac{A'}Q\right)\epsilon_Q\sqrt Q\,
e^{2\pi ir/q},\qquad
\epsilon_Q=\begin{cases}1,&Q\equiv1\pmod4,\\ i,&Q\equiv3\pmod4.
\end{cases}                                                    \tag{2}
\]
Here $(A'/Q)$ is the source-ring Jacobi symbol.  For
$\Gamma_k(A)=\sum_{x\bmod p^k}e^{2\pi iAx^2/p^k}$, split
$x=y+p^{k-1}z$.  The $z$-sum vanishes unless $p\mid y$ and gives
\[
\Gamma_k(A)=p\,\Gamma_{k-2}(A)\quad(k\ge2).
\]
Together with $\Gamma_1(A)=(A/p)\epsilon_p\sqrt p$ and $\Gamma_0=1$, this
proves every odd prime-power case.  CRT combines the prime-power sums, and
quadratic reciprocity reduces the accumulated signs to
$(A'/Q)\epsilon_Q$.  Thus (2) is an exact finite-sum identity.  If $Q=1$,
the separate constant branch is
$G_{q,n}=q e^{2\pi iC_n/q}$.  Thus (2), including its vanishing gate, holds for
every $q,a,b,n$ in the frozen family; finite sweeps are only regressions.

\section{Zero levels and a concrete specialization}
For an odd prime $p$, reduce $(A_n,B_n,C_n)$ modulo $p$.  If $A_n\ne0$,
completion of the square yields
\[
\#\{x:S_n\phi(x)=0\}=1+\left(\frac{\Delta}{p}\right),\qquad
\Delta=B_n^2-4A_nC_n.                                        \tag{3}
\]
If $A_n=0\ne B_n$ the count is one; if $A_n=B_n=0$, it is $p$ or zero
according as $C_n=0$ or not.  In the pure quadratic case $a=1,b=0$, $p\ge5$,
and $n\ne0\pmod p$,
\[
\Delta=\frac{n^2(1-n^2)}3,
\qquad \#\{x:S_nx^2=0\}=1+\left(\frac{\Delta}{p}\right).     \tag{4}
\]
Thus $n\equiv\pm1\pmod p$ gives the single double root, whereas
$n\equiv0\pmod p$ gives all $p$ roots.

\section{Same-clock unitary and boundary}
On $\mathcal H_q=\ell^2(\mathbb Z/q\mathbb Z)$ let
$(K_qf)(x)=f(x+1)$, let $M_\phi$ multiply by $e^{2\pi i\phi(x)/q}$, and put
$U_\phi=M_\phi K_q$.  Direct iteration gives the exact same-clock identity
\[
G_{q,n}(a,b)=\operatorname{Tr}(U_\phi^nK_q^{-n}).              \tag{5}
\]
The compensating shift is essential; (5) is not replaced by
$\operatorname{Tr}(U_\phi^n)$.
Let $P f(x)=f(-x)$, let $J$ be complex conjugation, set
$g(x)=(a-b)x^2$, and let $D_g$ multiply by $e^{2\pi ig(x)/q}$.  The
antiunitary $\Theta=D_gPJ$ obeys
\[
\Theta^2=I,\qquad \Theta U_\phi\Theta^{-1}=U_\phi^{-1}.         \tag{6}
\]
Indeed $g$ is even and
$g(x)-g(x-1)=\phi(-x)-\phi(x-1)=(a-b)(2x-1)$, exactly the multiplier
identity obtained when $P$ reverses the shift.  More explicitly,
\[
(\Theta U_\phi\Theta f)(x)
=e^{2\pi i\{g(x)-\phi(-x)-g(1-x)\}/q}f(x-1)
=e^{-2\pi i\phi(x-1)/q}f(x-1)=U_\phi^{-1}f(x).
\]

The strict tuple is $(\texttt{A1\_WEAK},\texttt{A2\_FAIL},
\texttt{A3\_FAIL},\texttt{A4\_NATURAL\_QUANTIZATION})$.
We claim no target trace/divisor/counting law, isolated stability determinant,
arithmetic local/Euler factor, root number, automorphy, Hilbert--P\'olya
construction, or Route-B authorization.  Scope:
\texttt{NO\_BAD\_EULER\_OR\_ROOT\_NUMBER}.

\section*{Declarations}
\textbf{Data availability.} All claim-bearing data and programs are packaged
with this certificate; no external dataset is used.  \textbf{Ethics.} No human
participants, animals, or sensitive personal data are involved.
\textbf{Author contributions/CRediT.} Anonymous technical certificate;
Conceptualization, Formal analysis, Software, Validation, and Writing are
recorded at package level.  \textbf{Competing interests.} None known.
\textbf{Funding.} No external funding is reported.  \textbf{AI-use
disclosure.} An AI coding assistant supported derivation planning, drafting,
and code review.  Packaged independent checks validate quantitative claims;
the assistant was not treated as an external peer reviewer.
\end{document}
